\documentclass[12pt,oneside]{book}  % 글씨 크기 13pt

% 언어 및 입력
\usepackage[english]{babel}
\usepackage[T1]{fontenc}
\usepackage[utf8]{inputenc}

% 가독성 좋은 단행본용 서체 (Libertinus 또는 Palatino)
\usepackage{libertinus}  % 또는: \usepackage[sc]{mathpazo}

\linespread{1.25}  % 줄간격: 단행본용

% 여백 설정 (A4, 단행본 스타일)
\usepackage[b5paper, margin=1in]{geometry}

% 수학 및 기호
\usepackage{amsmath, amssymb, amsthm, mathtools}

% 표와 배열
\usepackage{booktabs, tabularx}

% 인용 및 각주
\usepackage[notes, backend=biber]{biblatex-chicago}
\addbibresource{attention.bib}
\usepackage[flushmargin, hang]{footmisc}

% 기타 기능
\usepackage{csquotes}
\usepackage{enumitem}
\usepackage[hidelinks]{hyperref}

\usepackage{etoolbox}
\pretocmd{\section}{\clearpage}{}{}

% Footnotes reset only at each chapter input below.
% Do not reset by section/subsection: the Postscript must run 1--75.

% 헤더와 푸터
\usepackage{fancyhdr}
\pagestyle{fancy}
\fancyhf{}
\fancyfoot[C]{\thepage}

% 자동 번호매기기 제거: 섹션/장 번호 숨기기
\makeatletter
\renewcommand{\@seccntformat}[1]{}
\makeatother

% 섹션 번호링 없음: chapter, section, subsection 모두 제목만 표시
\setcounter{secnumdepth}{0}
\setcounter{tocdepth}{2}

% 문서 시작
\title{The Concept of Law, Third Edition}
\author{H. L. A. Hart}
\date{}

\begin{document}

\begin{titlepage}
  \centering
  \vspace*{3cm}
  {\Huge\textsc{The Concept of Law}}\\[1.5ex]
  {\Large Third Edition}\\[4ex]
  \textsc{H. L. A. Hart}\\[6ex]
  \vfill
\end{titlepage}

% ==========================
% 📚 본문 시작
% ==========================
\begin{document}

\frontmatter
\maketitle
\input{sections/01_1961_PREFACE.tex}
\input{sections/02_1994_Editors_Note.tex}
\input{sections/03_2012_Third_Preface.tex}
\setcounter{footnote}{0}
\input{sections/04_2012_INTRODUCTION.tex}

\tableofcontents

\clearpage

\mainmatter
\cleardoublepage
\markboth{}{}

\setcounter{footnote}{0}
\input{sections/chapter_1copy.tex}
\setcounter{footnote}{0}
\input{sections/chapter_2.tex}
\setcounter{footnote}{0}
\input{sections/chapter_3.tex}
\setcounter{footnote}{0}
\input{sections/chapter_4.tex}
\setcounter{footnote}{0}
\input{sections/chapter_5.tex}
\setcounter{footnote}{0}
\input{sections/chapter_6.tex}
\setcounter{footnote}{0}
\input{sections/chapter_7.tex}
\setcounter{footnote}{0}
\input{sections/chapter_8.tex}
\setcounter{footnote}{0}
\input{sections/chapter_9.tex}
\setcounter{footnote}{0}
\input{sections/chapter_10.tex}
\setcounter{footnote}{0}
\input{sections/chapter_11-POSTCSRIPT.tex}

\clearpage
\backmatter

\end{document}
